\documentclass[11pt]{scrartcl}

\usepackage[stixtwo]{fontsetup}
\usepackage{amsmath}

\begin{document}
\begin{enumerate}
	\item \[
		      a_n
		      = \frac{\left(\sqrt{n + \sqrt{n}} - \sqrt{n}\right)\left(\sqrt{n + \sqrt{n}} + \sqrt{n}\right)}{\sqrt{n + \sqrt{n}} + \sqrt{n}}
		      = \frac{\sqrt{n}}{\sqrt{n + \sqrt{n}} + \sqrt{n}}
		      = \frac{1}{\sqrt{1 + \frac{1}{\sqrt{n}}} + 1}
		      \to \frac{1}{2}
	      \]

	\item \[
		      b_n
		      = \frac{-\frac{1}{4} \cdot 16^n - 3 \cdot (-27)^n}{25^n}
		      = -\frac{1}{4} \left(\frac{16}{25}\right)^n
		      - 3 (-1)^n \left(\frac{27}{25}\right)^n
	      \]
	      Az első tag $0$-hoz tart, a második abszolút értéke $\infty$-hez, és
	      előjele váltakozik: páros $n$-re $b_n \to -\infty$, páratlan $n$-re
	      $b_n \to \infty$. A sorozatnak tehát két torlódási pontja van
	      ($\pm \infty$), így nincs határértéke.
\end{enumerate}
\end{document}
